\chapter{Calprotectin (Stool)}
\section*{What Is This Test?}
Fecal calprotectin is a noninvasive biomarker that measures the concentration of calprotectin, a 36.5-kDa calcium- and zinc-binding protein, in stool. Calprotectin constitutes approximately 60\% of the cytosolic protein content of neutrophils and is also expressed in monocytes, activated macrophages, and mucosal epithelial cells. It is released from neutrophils during intestinal mucosal inflammation and is remarkably resistant to degradation by intestinal and bacterial proteases, remaining stable in stool for up to 7 days at room temperature. Fecal calprotectin is used primarily to differentiate irritable bowel syndrome (IBS, a functional gastrointestinal disorder with no intestinal inflammation --- normal fecal calprotectin) from inflammatory bowel disease (IBD, Crohn disease, and ulcerative colitis --- elevated fecal calprotectin) in patients presenting with chronic abdominal pain, diarrhea, or altered bowel habits. It is also used to monitor disease activity in established IBD and to predict clinical relapse. Fecal calprotectin correlates well with endoscopic and histologic evidence of intestinal inflammation (especially colonic inflammation; it is less sensitive for isolated small bowel Crohn disease). It has largely replaced the need for invasive colonoscopy to exclude IBD when the pre-test probability is low, and a normal fecal calprotectin (<50 mcg/g) effectively excludes IBD with a high negative predictive value.
\section*{Alternative Names}
Fecal Calprotectin; Stool Calprotectin; Fecal Calprotectin Assay; IBD Screen
\section*{Principle}
Fecal calprotectin is measured from a single random stool sample (1--5 g) by quantitative ELISA (enzyme-linked immunosorbent assay) or automated chemiluminescent immunoassay (CLIA) using monoclonal antibodies against calprotectin. Results reported in mcg calprotectin/g stool. Point-of-care lateral flow tests (semi-quantitative) are also available.
\section*{History}
Calprotectin was discovered in 1980 by Fagerhol and colleagues, who isolated the dominant cytosolic protein of neutrophils (originally termed the L1 protein; now recognized as the S100A8/S100A9 heterodimer, also known as MRP-8/14). In 1997, Roseth and colleagues demonstrated that calprotectin resists proteolytic degradation and can be quantified in stool, and that levels correlate with intestinal inflammation. Studies by Tibble and colleagues in the early 2000s established that fecal calprotectin reliably distinguishes organic inflammatory bowel disease (IBD) from functional irritable bowel syndrome (IBS), reducing the need for colonoscopy. The 2013 NICE diagnostic guidance (DG11) endorsed fecal calprotectin testing in primary care to determine which patients warrant referral for endoscopy. Fecal calprotectin is now recommended in European and American IBD guidelines as a noninvasive marker of mucosal inflammation.
\section*{Why This Test Is Ordered}
\begin{itemize}
  \item Differentiate IBD (Crohn disease, ulcerative colitis) from IBS in patients with chronic abdominal pain, diarrhea, or altered bowel habits, thereby avoiding unnecessary colonoscopy when the level is normal
  \item Estimate the probability of IBD in primary care before referral to a gastroenterologist
  \item Monitor disease activity in patients with established IBD and assess response to medical therapy (e.g., biologics, aminosalicylates)
  \item Predict clinical relapse in IBD: a rising fecal calprotectin precedes endoscopic and clinical relapse by 3--6 months
  \item Evaluate chronic diarrhea or suspected intestinal inflammation in children
  \item Help select which patients with unexplained gastrointestinal symptoms require urgent versus elective endoscopy
\end{itemize}
\section*{Is This a Primary or Secondary Test}
This is a primary screening test used to triage patients with lower gastrointestinal symptoms before endoscopy, and a secondary monitoring test in patients with established IBD. It is not diagnostic of IBD itself; colonoscopy with biopsies remains the gold standard.
\section*{How This Test Is Performed}
A single random stool sample (1--5 g) is collected in a plain screw-top container; no preservative is required because calprotectin is stable in stool for up to 7 days at room temperature. The sample is transported to the laboratory, extracted, and analyzed by quantitative ELISA or automated chemiluminescent immunoassay (CLIA) using monoclonal antibodies against calprotectin. The result is reported in micrograms of calprotectin per gram of stool (mcg/g). Batch immunoassay results are typically available within 1--3 days, while point-of-care lateral flow devices provide semi-quantitative results within minutes.
\section*{What Preparation Is Needed}
No fasting or special diet is required, and the sample may be collected at any time of day. If possible, nonsteroidal anti-inflammatory drugs (NSAIDs) and aspirin should be avoided for several days before collection, because NSAID enteropathy can elevate calprotectin. The sample should be kept cool or at room temperature (not frozen) before transport.
\section*{Contraindications and Precautions}
\begin{itemize}
  \item Active gastrointestinal bleeding or fresh blood in stool can falsely elevate fecal calprotectin
  \item NSAIDs, aspirin, and proton pump inhibitors may produce mild elevations
  \item Children under 4 years have higher normal values than older children and adults
  \item Acute infectious gastroenteritis can transiently raise calprotectin; defer testing until resolution
  \item Sensitivity is reduced for isolated small bowel Crohn disease; a normal value does not fully exclude small bowel inflammation
  \item Results are not useful for determining disease location or extent
\end{itemize}
\section*{Risks}
None. Fecal calprotectin testing is entirely noninvasive and carries no physical risk; the only inconvenience is collection and handling of a stool specimen.
\section*{What Happens After the Test}
A normal result ($<$50 mcg/g) effectively excludes IBD and avoids colonoscopy, whereas a borderline or elevated result prompts endoscopic evaluation or repeat testing. In patients with established IBD, serial measurements guide therapy escalation and predict relapse. Repeat fecal calprotectin is typically measured every 3--6 months during disease monitoring.
\section*{Understanding Results}
\subsection*{Normal (<50 mcg/g)}
Minimal or no neutrophilic intestinal inflammation. IBS is likely; IBD is unlikely (NPV >95\% for IBD). A normal fecal calprotectin level in a patient with chronic GI symptoms effectively excludes IBD and avoids unnecessary colonoscopy.
\subsection*{Borderline (50--120 mcg/g)}
Mild intestinal inflammation. May represent mild/treated/quiescent IBD, NSAID enteropathy, microscopic colitis, infectious gastroenteritis (bacterial), colorectal adenomas/carcinoma, or diverticulitis. Repeat testing and clinical correlation are indicated.
\subsection*{Elevated (120--250 mcg/g)}
Active intestinal inflammation. Consistent with IBD or other organic GI pathology. Colonoscopy with biopsies is typically warranted.
\subsection*{Markedly Elevated (>250 mcg/g)}
Significant neutrophilic intestinal inflammation. Highly suggestive of active IBD. Colonoscopy with biopsies is indicated. Fecal calprotectin >250 mcg/g in a patient with known IBD indicates active mucosal inflammation and a high risk of clinical relapse within the next 3--6 months. Escalation of IBD therapy should be considered.
\section*{Normal Results}
Normal: <50 mcg/g (negative). Borderline: 50--120 mcg/g. Elevated (IBD likely): >120--250 mcg/g. Markedly elevated: >250 mcg/g. Values are age-dependent: children <4 years have higher normal fecal calprotectin levels. NSAIDs and aspirin can elevate fecal calprotectin (NSAID enteropathy). Proton pump inhibitors (PPIs) cause mild calprotectin elevation.
\section*{Follow-up Tests}
Colonoscopy with ileoscopy and biopsies (gold standard for IBD diagnosis), upper endoscopy (for suspected Crohn disease with upper GI involvement), CT or MR enterography (for small bowel evaluation), fecal lactoferrin (alternative fecal inflammatory marker), serum CRP and ESR, and repeat fecal calprotectin in 3--6 months (for IBD monitoring and to predict relapse).
\section*{References}
\begin{enumerate}
  \item Tibble JA, Sigthorsson G, Bridger S, et al. Surrogate markers of intestinal inflammation are predictive of relapse in patients with inflammatory bowel disease. Gastroenterology. 2000;119(1):15--22.
  \item Waugh N, Cummins E, Royle P, et al. Faecal calprotectin testing for differentiating amongst inflammatory and non-inflammatory bowel diseases. Health Technol Assess. 2013;17(55):1--211.
  \item D'Haens G, Ferrante M, Vermeire S, et al. Fecal calprotectin is a surrogate marker for endoscopic lesions in inflammatory bowel disease. Inflamm Bowel Dis. 2012;18(12):2218--2224.
  \item van Rheenen PF, Van de Vijver E, Fidler V. Faecal calprotectin for screening of patients with suspected inflammatory bowel disease: diagnostic meta-analysis. BMJ. 2010;341:c3369.
  \item Roseth AG, Aadland E, Jahnsen J, Raknerud N. Assessment of disease activity in ulcerative colitis by faecal calprotectin, a novel granulocyte marker protein. Digestion. 1997;58(2):176--180.
  \item Fagerhol MK. Calprotectin, a faecal marker of organic gastrointestinal abnormality. Lancet. 2000;356(9244):1783--1784.
\end{enumerate}
\clearpage
